%%%%%%%%%%%%%%%%%%%%%%%%%%%%%%%%%%%%%%%%%%%%%%%%%%%%%%%%%%%%%%%%%%%%%%%%%%%%%%%
\section{Threat Model}\label{sec:threats}
%%%%%%%%%%%%%%%%%%%%%%%%%%%%%%%%%%%%%%%%%%%%%%%%%%%%%%%%%%%%%%%%%%%%%%%%%%%%%%%

%The original threat model is the next 3 paragraphs
%We describe KBID as an authentication mechanism that requires a wearable device that transmits short-range authentication data via BCC. We recognize confidentiality, integrity, and availability as security goals specific to KBID. Specifically, data stored on the device and transmitted to the authentication module should be kept secret from and not modifiable by unauthorized entities, and the data should be accessible to both bracelet and authentication module. We omit the privacy of the authenticating user because the system must validate her access to the system or resource. 

%Adversaries are typically distinguished based on their goals, capabilities, and relation to a system. We define the following adversarial classification criteria for KBID: active adversaries that are able to read, modify, and inject communication between the device and authentication module, and passive adversaries that can eavesdrop on the communication; internal entities that have legitimate access to the device, and external entities that do not; single or coordinated group entities, and; sophisticated adversaries with specialized equipment (e.g., high-gain directional antenna or unauthorized authentication module), and unsophisticated adversaries with common equipment.

%The KBID device or authentication module may both be used as attack surfaces. For example, an adversary may disrupt KBID authentication by physically damaging the wearable device or authenticator module. We classify this and other KBID security threats into the following categories: BCC threats, whereby the adversary is able to passively eavesdrop on communication, or actively jam, replay, modify, forge, or drop communication; hardware threats, whereby the adversary is able to induce incorrect outputs from a valid device, and; software threats, whereby the adversary is able to alter the logic of KBID's device, authentication module, or client software through software vulnerabilities.

%New Threat Model
As a hardware and software solution, the threat model for KBID includes many subjects that apply to any such system.  These include attack types (denial of service, message forging or tampering, hardware tampering, and others) as well as a study of potential adversaries and other topics.  Here we focus on two threats that are unique to KBID.  These threats require an active adversary that has the ability to get close to where KBID is used.  It is possible to impersonate an authentication module.  An attacker could use a counterfeit authentication module that issues \emph{Get Status} commands when a user touches something connected to it, e.g. a door knob.  We discuss mitigating this vector below. Second, while we are using body coupled communication to transmit data without emitting RF, it could be the case that the user's body acts as a broadcast antenna and emits the data into the environment where it could conceivably be intercepted.  We plan on assessing the reality of this treat during the development of the next prototype.
